\documentclass[11pt,a4paper]{article}

\usepackage[utf8]{inputenc}
\usepackage[T1]{fontenc}
\usepackage[polish]{babel}
\usepackage[margin=2.5cm]{geometry}
\usepackage{enumitem}

% oznaczenia autorstwa: [IK] Igor Kosakowski, [HW] współautor (sprzęt), [wsp.] wspólne
\newcommand{\IK}{\textsuperscript{\textbf{[IK]}}}
\newcommand{\HW}{\textsuperscript{\textbf{[HW]}}}
\newcommand{\WS}{\textsuperscript{\textbf{[wsp.]}}}

% lista numerowana 1. / 1.1.
\newlist{toc}{enumerate}{2}
\setlist[toc,1]{label=\textbf{\arabic*.},leftmargin=2em,itemsep=6pt}
\setlist[toc,2]{label=\arabic{toci}.\arabic*.,leftmargin=2.4em,itemsep=0pt,topsep=2pt}

\setlength{\parindent}{0pt}
\pagestyle{empty}

\begin{document}

\section*{Proponowany spis treści}

\begin{toc}
\item \textbf{Wstęp} \WS
  \begin{toc}
    \item Motywacja
    \item Cel i zakres pracy
    \item Podział pracy między autorów
  \end{toc}

\item \textbf{Podstawy teoretyczne} \WS
  \begin{toc}
    \item Zjawisko \emph{rolling shutter}
    \item Model kamery i geometria obrazowania
    \item Uczenie maszynowe w przetwarzaniu obrazu
  \end{toc}

\item \textbf{Przegląd literatury} \WS
  \begin{toc}
    \item Metody korekcji efektu \emph{rolling shutter}
    \item Zbiory danych
  \end{toc}

\item \textbf{Stanowisko pomiarowe i dane rzeczywiste} \HW
  \begin{toc}
    \item Konstrukcja stanowiska i akwizycja obrazu
    \item Kalibracja i synchronizacja kamer
    \item Zbiór danych rzeczywistych
  \end{toc}

\item \textbf{Generowanie danych syntetycznych} \IK
  \begin{toc}
    \item Założenia i model zjawiska
    \item Algorytm generowania danych
    \item Opis wygenerowanych zbiorów
  \end{toc}

\item \textbf{Trening i porównanie modeli} \IK
  \begin{toc}
    \item Rozważane architektury
    \item Proces trenowania
    \item Napotkane problemy i wprowadzone usprawnienia
  \end{toc}

\item \textbf{Ewaluacja i wyniki} \WS
  \begin{toc}
    \item Metodyka i metryki oceny
    \item Wyniki na danych syntetycznych
    \item Wyniki na danych rzeczywistych
    \item Dyskusja
  \end{toc}

\item \textbf{Podsumowanie i kierunki dalszych prac} \WS

\item[] \textbf{Bibliografia}
\item[] \textbf{Dodatki}
\end{toc}

\end{document}
